\section{Distributed Key Generation (DKG)}
\label{dkg-reqs}

Informally, a \emph{Distributed Key Generation} (DKG) scheme allows a set of $n$ players to jointly generate a public key $\pk$ and a secret key $\sk$,
such that all parties learn the public key $\pk$ but no single party learns $\sk$.
Instead, each party $\id_i$ has a share $\sk_i$ of the secret key,
such that a threshold $t$ number of shares are required to recover $\sk$,
where $t \leq n$.
A DKG generates secret and public keys that are compatible with some \emph{target} key generation scheme which operates in the centralized (single-party) setting.

We now formalize the definition of a DKG.
We begin by defining a target single-party key generation scheme $\targetkeygen$ that the DKG implements in a distributed manner.

\begin{definition}
A \defn{target key generation scheme} $\targetkeygen$ is the tuple $\targetkeygen = (\keygen, \verify)$, where:
\begin{itemize}
    \item $\keygen(\lambda) \rightarrow (\SK, \PK)$: A probabilistic algorithm that
      accepts as input a security parameter $\lambda$ and outputs a secret and public keypair $(\SK, \PK)$, where $\SK$ is an element of a secret domain $\secretspace$ and $\PK$ is an element of the public domain $\publicspace$.
    \item $\verify(\SK, \PK) \rightarrow \zo$: A deterministic algorithm that accepts as input a secret and public keypair $(\SK, \PK)$,
  and outputs $0$ if the keypair is invalid; otherwise, outputs~$1$.
    \item For correctness, we require that for all $(\SK, \PK) \gets \targetkeygen.\keygen(\lambda)$, we have
      \[
        \targetkeygen.\verify(\SK, \PK) = 1
      \]
\end{itemize}
\end{definition}



\begin{definition}
  A \defn{distributed key generation scheme}, or \defn{DKG},
  is the tuple of the protocols $\dkg = (\keygen, \recover)$
  parameterized by the number of rounds $\numrounds$.
  A DKG is defined with respect to a target key generation scheme $\targetkeygen$,
  %and initialized with the hash function $\Hash$, whose domain and codomain depend on $\dkg$,
  where:

  \begin{itemize}
    \item $\keygen(\lambda, n, t) \rightarrow  \langle (\PK, \qual, \auxinfo), (\status_i,\sk_i)_{i \in \set{n}}\rangle$:
      An interactive protocol between $n$ participants.
      Accepts as input the security parameter $\lambda$,
      and positive integers $n \geq t \geq 1$,
      where $n$ is the number of participants and $t$ is the threshold.

      The protocol outputs
      a public value $(\PK, \qual, \auxinfo)$
      and private values $(\status_i, \sk_i), i \in \set{n}$
      where:
      \begin{itemize}
        \item $\PK$ is the public key representing the group;
          it will be $\bot$ if the protocol terminated unsuccessfully.
        \item $\qual \subseteq \set{n}$ is the set of qualified players remaining at the end of the protocol;
          it will be $\bot$ if the protocol terminated unsuccessfully.
        \item $\auxinfo$ is any auxiliary information that may be required for a particular application,
          such as participant public keys,
        \item $\status_i$ indicates if participant $i$ completed the protocol successfully,
      and can be either $\abort$, $\fail$, or $\accept$.
      Here, $\abort$ indicates the protocol terminated early,
      $\fail$ indicates a critical fail state has been reached,
      and $\accept$ indicates that the protocol completed successfully.
        \item $\SK_i$ is participant $i$'s secret key share;
          it will be $\bot$ if the protocol terminated unsuccessfully.
      \end{itemize}


      % CK- I originally had this, but this should be defined in the construction, not the definition.
      %Note that once a participant's state has been set to $\fail$ or $\abort$,
      %the participant no longer performs subsequent round operations for that session (i.e, its protocol execution is terminated).
      %\ian{I don't see anywhere in the paper where $\status_i$ can be set to $\fail$? If there is such, be more clear about the difference between that and $\abort$.}
      % CK- we don't have that in our DKG, but others could.
      % CK- I said that abort indicates the protocol terminated early and fail is a critical failure.
      %\ian{This is a per-player output, though; if the protocol completes, but say 3 players are kicked out, what do the kicked-out players get as their private output?  What do the non-kicked-out players get?}
      % CK- We don't define what kicked out players do, because at this point, they are considered adversarial.
      % CK- we do define what honest players statuses get set to in the construction.
      % The definition requires just defining the inputs/outputs. The logic for how these get set is defined in the constrcution.


      \medskip
      Internally to the $\keygen$ protocol, each participant performs the following:

      \begin{itemize}
        \item $\dkgkeygen_0(\lambda, n, t, \id_i) \rightarrow (\party_i,\outvecp,\bmsg_i)$:
          The algorithm run by the participant to initiate the protocol.
          Accepts the security parameter $\lambda$,
          the number of participants $n$, the threshold $t$, and the participant identifier $\id_i$.
          Outputs the state $\party_i$ for that participant,
          the vector $\outvecp$ of any outgoing peer-to-peer messages to other parties,
          and a broadcast message $\bmsg_i$ to all other parties.
          See Section~\ref{prelim:network} for definitions of these network channels.
	  %\ian{The vector $\outvecp$ can also be $\bot$, right? It sometimes is in the actual protocol.}
    % CK- yes, but so can any variable, we don't need to say that explicitly here.
        \item $\dkgkeygen_k(\party_i, \invecp, \invecb) \rightarrow (\party_i, \outvecp, \bmsg_i)$:
          The algorithm run for each intermediate round $k \in \{ 1, \ldots, \numrounds-1\}$ in the protocol.
          Accepts as input the participant's internal state $\party_i$,
          the vector $\outvecp$ of any peer-to-peer messages sent by other participants in the previous round,
          and the set $\invecb$ of any broadcast messages sent by other participants in the previous round.
          Outputs the tuple $(\party_i, \outvecp, \bmsg_i)$.
        \item $\finalize(\party_i, \invecp, \invecb) \rightarrow \langle(\PK, \qual,\auxinfo), (\status_i, \sk_i)_{i \in \set{n}}\rangle$:
          Accepts participant state $\party_i$, peer-to-peer messages $\outvecp$, and broadcast messages $\invecb$.
          Outputs a public output $(\PK, \qual,\auxinfo)$ and a private output $(\status_i, \SK_i)$.
      \end{itemize}

    \item $\recover(\PK, t, \{ (i, \SK_i) \}_{i \in \coalition}) \rightarrow \SK/\fail$:
      The algorithm run by the participant to finish its protocol execution.
      Accepts as input the public key $\PK$,
      the threshold $t$,
      and a set of secret key shares $\{(i,\sk_i) \}_{i \in C}$,
      where $\lvert C \rvert \geq t$.
      Derive $\SK$ from $\{\SK_i \}_{i \in \coalition}$.
      If $\targetkeygen.\verify(\SK, \PK) \neq 1$, output $\fail$.
      Otherwise, output $\SK$.
    \item For correctness, we require that, when all participants follow the protocol:
\begin{align*}
    \text{if } & \dkg.\keygen(\lambda, n, t) \rightarrow \left((\PK, \qual), \langle (\status_i, \sk_i)\rangle_{i \in \set{n}}\right) \\[1mm]
    \text{then } & \dkg.\recover(\PK, t, \{(i,\SK_i)\}_{i \in \coalition}) \rightarrow \SK \text{ for all } \coalition \subseteq \set{n}, \lvert \coalition \rvert \geq t, \\[1mm]
    \text{and } &  \targetkeygen.\verify(\SK, \PK) \rightarrow 1
 \end{align*}
\end{itemize}
\end{definition}

A DKG must be correct, and fulfill the security notions of
zero-knowledge and indistinguishability.
Additionally, a DKG may fulfill the notion of either strong or weak robustness.
We explain these concepts next.

